\documentclass[12pt, a4paper]{article}

\usepackage[utf8]{inputenc}
\usepackage[spanish]{babel}
\usepackage{amsmath, amssymb, braket, physics}
\usepackage{geometry}
\geometry{a4paper, margin=2.5cm}
\usepackage{titlesec}
\usepackage{enumitem}
\usepackage{listings}
\usepackage{xcolor}

% Configuración de estilo de código
\definecolor{codegreen}{rgb}{0,0.6,0}
\definecolor{codegray}{rgb}{0.5,0.5,0.5}
\definecolor{codepurple}{rgb}{0.58,0,0.82}
\definecolor{backcolour}{rgb}{0.95,0.95,0.92}

\lstdefinestyle{mystyle}{
    backgroundcolor=\color{backcolour},   
    commentstyle=\color{codegreen},
    keywordstyle=\color{magenta},
    numberstyle=\tiny\color{codegray},
    stringstyle=\color{codepurple},
    basicstyle=\ttfamily\footnotesize,
    breakatwhitespace=false,         
    breaklines=true,                 
    captionpos=b,                    
    keepspaces=true,                 
    numbers=left,                    
    numbersep=5pt,                  
    showspaces=false,                
    showstringspaces=false,
    showtabs=false,                  
    tabsize=2
}
\lstset{style=mystyle}

% Configuración de secciones
\titleformat{\section}{\large\bfseries}{\thesection}{1em}{}
\titleformat{\subsection}{\normalsize\bfseries}{\thesubsection}{1em}{}
\setlength{\parindent}{0pt}
\setlength{\parskip}{0.5em}

\title{\textbf{Syllabus: Conceptos Básicos de Información Cuántica}}
\author{Basado en: \textit{Basics of Quantum Information} (IBM Quantum) \\ Nivel: Preparatoria Avanzada}
\date{\today}

\begin{document}

\maketitle

\noindent\textbf{Duración Total:} 2 Horas (1 hr Teoría + 1 hr Laboratorio Práctico) \\
\textbf{Requisitos:} Conocimientos previos en álgebra elemental y nociones de programación en Python.

\section*{Resumen y Enfoque del Curso}
Este curso intensivo adapta los fundamentos del programa de IBM Quantum para estudiantes de preparatoria. La primera mitad proporciona un marco teórico y físico, partiendo de la dualidad onda-partícula y la ecuación de Schrödinger, para luego formalizar matemáticamente los conceptos de probabilidad, números complejos y matrices aplicados a sistemas de dos niveles (Qubits). 

La segunda mitad es un laboratorio práctico en Python. Los estudiantes traducirán los conceptos matriciales a código utilizando arreglos (NumPy) y posteriormente construirán, simularán y medirán sus primeros circuitos cuánticos utilizando el framework Qiskit de IBM. Se mantendrá en todo momento un enfoque que favorezca el razonamiento matemático detallado y de alto nivel.

\vspace{0.5cm}

\hrule
\vspace{0.5cm}

\section*{Temario: Parte I - Fundamentos Físicos y Matemáticos (Hora 1)}

\subsection*{1.1 De la Física Clásica a la Cuántica (15 min)}
\begin{itemize}
    \item \textbf{Dualidad Onda-Partícula:} El experimento de la doble rendija y el colapso de la función de onda inducido por la observación.
    \item \textbf{La Ecuación de Schrödinger:} Evolución temporal determinista de los estados cuánticos expresada como $i\hbar \frac{\partial}{\partial t}\ket{\Psi(t)} = \hat{H}\ket{\Psi(t)}$.
    \item \textbf{Probabilidad Cuántica:} Contraste entre la adición de probabilidades clásicas ($P = p_1 + p_2$) y la interferencia de amplitudes de probabilidad cuánticas ($P = |\psi_1 + \psi_2|^2$).
\end{itemize}

\subsection*{1.2 Herramientas Matemáticas: Complejos y Matrices (20 min)}
\begin{itemize}
    \item \textbf{Números Complejos:} Representación rectangular ($z = x + iy$) y polar ($z = r e^{i\theta}$). La magnitud al cuadrado como métrica de probabilidad: $|z|^2 = z^* z$.
    \item \textbf{Álgebra Lineal para Cuántica:} Vectores columna (Kets) y vectores fila conjugados (Bras). 
    \item \textbf{Matrices Unitarias:} Operadores que preservan la probabilidad, cumpliendo $U^\dagger U = I$.
\end{itemize}

\subsection*{1.3 El Qubit y la Información Cuántica (25 min)}
\begin{itemize}
    \item \textbf{Definición del Qubit:} Combinación lineal de estados base $\ket{\psi} = \alpha\ket{0} + \beta\ket{1}$, donde $\alpha, \beta \in \mathbb{C}$ y $|\alpha|^2 + |\beta|^2 = 1$.
    \item \textbf{Compuertas de un Qubit:} Representación matricial de las compuertas $X$ (NOT cuántico) y Hadamard ($H$).
    \item \textbf{Sistemas Multiqubit y Entrelazamiento:} Introducción rápida al producto tensorial y los estados de Bell.
\end{itemize}

\newpage

\section*{Temario: Parte II - Laboratorio Práctico en Python (Hora 2)}

\subsection*{2.1 Álgebra Lineal Computacional con NumPy (20 min)}
\begin{itemize}
    \item \textbf{Vectores y Matrices en Python:} Definición manual de los estados $\ket{0}$ y $\ket{1}$ como arreglos de NumPy.
    \item \textbf{Operaciones Matriciales:} Aplicación del producto punto matricial (\texttt{np.dot} o \texttt{@}) para simular la acción de la compuerta $X$ y $H$ sobre el estado inicial.
\end{itemize}
\begin{lstlisting}[language=Python, caption={Ejemplo ilustrativo para el manual}]
import numpy as np
ket_0 = np.array([1, 0])
H = (1/np.sqrt(2)) * np.array([[1, 1], [1, -1]])
estado_superposicion = H @ ket_0
\end{lstlisting}

\subsection*{2.2 Introducción a IBM Qiskit (20 min)}
\begin{itemize}
    \item \textbf{El Objeto QuantumCircuit:} Inicialización de registros cuánticos y clásicos (\texttt{QuantumCircuit(q, c)}).
    \item \textbf{Traducción de Compuertas:} Mapeo de la teoría matemática a funciones de la librería (\texttt{qc.h(0)}, \texttt{qc.cx(0,1)}).
    \item \textbf{Construcción de un Estado Entrelazado:} Creación paso a paso del circuito para el estado de Bell $\ket{\Phi^+} = \frac{1}{\sqrt{2}}(\ket{00} + \ket{11})$.
\end{itemize}

\subsection*{2.3 Simulación, Medición y Visualización (20 min)}
\begin{itemize}
    \item \textbf{Medición (Regla de Born en código):} Colapso del estado cuántico y almacenamiento de la información en bits clásicos usando \texttt{qc.measure()}.
    \item \textbf{Ejecución en Simulador:} Uso del \texttt{AerSimulator} para ejecutar el circuito miles de veces (shots).
    \item \textbf{Análisis de Resultados:} Graficación de las distribuciones de probabilidad resultantes usando \texttt{plot\_histogram}. Demostración práctica de la aleatoriedad cuántica estructurada.
\end{itemize}

\end{document}
